% RSE Supplementary: the analyses of pre-registration Amendment 20, made in answer to the journal's
% reviewers (RSASE-D-26-03629; docs/DECISIONS.md D-161). Exploratory; specified before they were
% computed. Every number is a \metric{} macro (keys rev3.*, tools/metrics_rev_x3.py).

All analyses in this section answer the journal's reviewers. They were specified in Amendment 20 of the pre-registration before any of them was computed. They are exploratory, none replaces a pre-specified verdict, and the decision rule of the article (Eq.~\ref{M-eq:rule}) is unchanged. None uses new satellite processing: each reads data that the earlier analyses had already produced. Figure and section numbers refer to the main article.

\subsection*{The audit without the deltaic sectors}
\emph{Why.} In the Volga delta sector the water extent depends on the river's discharge as well as on the sea's level. The sector holds the delta plain, the lakes west of it and the lower Volga--Akhtuba floodplain upstream of Astrakhan. Water levels in the delta vary by metres with the river's flow and with wind surges, and flood large areas of the flat delta plain \citep{yaitskaya2023}. The Ural delta sector holds ground of the same kind on a smaller area. A waterline mapped on such ground is not a contour of the sea's level. Both sectors were labelled deltaic before any evaluation, and no claim of the article rests on them. Their cells nevertheless entered every domain-wide quantity: each grid's area-fitted offset is fitted on the training blocks of the whole domain, of whose known cells \metric{rev3.marine.fit_cells_deltaic_pct}\% lie in the two deltaic sectors, and the domain area bias, the observed 2026 sea and the exposure counts include them. We therefore repeated these quantities on the \emph{marine domain}: the study domain less the two deltaic sectors. The mask is the pre-specified sector polygons, so it involves no new choice.

\emph{Offsets and verdicts.} Refitted on the training blocks of the marine domain, the offsets change by \metric{rev3.marine.c_change_min_m} to \metric{rev3.marine.c_change_max_m}~m, and those of \metric{rev3.marine.n_grids_unchanged} grids do not change (Supplementary Table~\ref{tab:a20_marine}). One pass over the levels then repeated the audit at the refitted offsets with the original noise floor. Run at the original offsets, the same pass reproduces every per-level displacement and floor of the main analysis, every verdict, every sector median and every block median. At the refitted offsets none of the \metric{rev3.marine.n_cases_checked} grid-by-checked-sector verdicts changes. Of all \metric{rev3.marine.n_cases} verdicts, \metric{rev3.marine.n_changed} change, each by a single level and neither in a checked sector: ETOPO~2022 now \metric{rev3.marine.changed.etopo2022.ural} in the Ural delta (\metric{rev3.marine.changed.etopo2022.ural.n_exceeding} of \metric{rev3.marine.changed.etopo2022.ural.n_evaluable} levels above the floor, against \metric{rev3.marine.changed.etopo2022.ural.n_exceeding_primary}), and SRTM15+ now \metric{rev3.marine.changed.srtm15p.komsomolets} in Komsomolets Bay (\metric{rev3.marine.changed.srtm15p.komsomolets.n_exceeding} of \metric{rev3.marine.changed.srtm15p.komsomolets.n_evaluable}, against \metric{rev3.marine.changed.srtm15p.komsomolets.n_exceeding_primary}). The four failing grids' sector-median displacement becomes \metricdp{rev3.marine.headline.marine.west.d_km_min}{1}--\metricdp{rev3.marine.headline.marine.west.d_km_max}{1}~km on the Kizlyar--Dagestan coast (\metricdp{rev3.marine.headline.shape.west.d_km_min}{1}--\metricdp{rev3.marine.headline.shape.west.d_km_max}{1}~km in the main analysis), \metricdp{rev3.marine.headline.marine.interfluve.d_km_min}{1}--\metricdp{rev3.marine.headline.marine.interfluve.d_km_max}{1}~km on the interfluve coast (\metricdp{rev3.marine.headline.shape.interfluve.d_km_min}{1}--\metricdp{rev3.marine.headline.shape.interfluve.d_km_max}{1}~km) and \metricdp{rev3.marine.headline.marine.northeast.d_km_min}{1}--\metricdp{rev3.marine.headline.marine.northeast.d_km_max}{1}~km on the northeastern coast (\metricdp{rev3.marine.headline.shape.northeast.d_km_min}{1}--\metricdp{rev3.marine.headline.shape.northeast.d_km_max}{1}~km). These are ranges over the four grids. The largest move of a single grid is that of GEBCO\_2022 on the northeastern coast, from \metric{rev3.marine.four.coasts3.d_move_max_from_km} to \metric{rev3.marine.four.coasts3.d_move_max_to_km}~km for a change of \metric{rev3.marine.gebco2022.c_change_m}~m in its offset, or \metric{rev3.marine.four.coasts3.d_move_over_q_max} of that coast's noise floor; its verdict there does not change.

\begin{table}[!htbp]
\centering
\small
\setlength{\tabcolsep}{4pt}
\caption{\textbf{The audit on the marine domain} (Amendment 20 item 1): the study domain less the Volga delta and Ural delta sectors. $c$: area-fitted offset on the training blocks of the whole domain (main analysis) and of the marine domain. Sectors failed: sectors in which the grid fails the shoreline test (Eq.~\ref{M-eq:rule}) at each offset. Area bias: median over the level units of (predicted $-$ observed) sea as a percentage of the observed sea, as distributed, over the whole domain and over the marine domain; and, last column, over the marine domain at the marine-domain offset.}
\label{tab:a20_marine}
\begin{tabular}{lccccccccc}
\hline
 & \multicolumn{2}{c}{$c$ (m)} & \multicolumn{2}{c}{Failed, of seven} & \multicolumn{2}{c}{Failed, of four checked} & \multicolumn{3}{c}{Area bias (\%)} \\
Grid & main & marine & main & marine & main & marine & domain & marine & marine, corrected \\
\hline
GEBCO\_2026 & \metric{h1.shape_offset.gebco2026_m} & \metric{rev3.marine.gebco2026.c_marine_m} & \metric{rev3.marine.gebco2026.shape.n_all7} & \metric{rev3.marine.gebco2026.marine.n_all7} & \metric{rev3.marine.gebco2026.shape.n_checked4} & \metric{rev3.marine.gebco2026.marine.n_checked4} & \metric{rev3.marine.area_bias.domain.total0.gebco2026.median_pct} & \metric{rev3.marine.area_bias.marine.total0.gebco2026.median_pct} & \metric{rev3.marine.area_bias.marine.marine.gebco2026.median_pct} \\
GEBCO\_2022 & \metric{h1.shape_offset.gebco2022_m} & \metric{rev3.marine.gebco2022.c_marine_m} & \metric{rev3.marine.gebco2022.shape.n_all7} & \metric{rev3.marine.gebco2022.marine.n_all7} & \metric{rev3.marine.gebco2022.shape.n_checked4} & \metric{rev3.marine.gebco2022.marine.n_checked4} & \metric{rev3.marine.area_bias.domain.total0.gebco2022.median_pct} & \metric{rev3.marine.area_bias.marine.total0.gebco2022.median_pct} & \metric{rev3.marine.area_bias.marine.marine.gebco2022.median_pct} \\
ETOPO~2022 & \metric{h1.shape_offset.etopo2022_m} & \metric{rev3.marine.etopo2022.c_marine_m} & \metric{rev3.marine.etopo2022.shape.n_all7} & \metric{rev3.marine.etopo2022.marine.n_all7} & \metric{rev3.marine.etopo2022.shape.n_checked4} & \metric{rev3.marine.etopo2022.marine.n_checked4} & \metric{rev3.marine.area_bias.domain.total0.etopo2022.median_pct} & \metric{rev3.marine.area_bias.marine.total0.etopo2022.median_pct} & \metric{rev3.marine.area_bias.marine.marine.etopo2022.median_pct} \\
ETOPO1 & \metric{h1.shape_offset.etopo1_m} & \metric{rev3.marine.etopo1.c_marine_m} & \metric{rev3.marine.etopo1.shape.n_all7} & \metric{rev3.marine.etopo1.marine.n_all7} & \metric{rev3.marine.etopo1.shape.n_checked4} & \metric{rev3.marine.etopo1.marine.n_checked4} & \metric{rev3.marine.area_bias.domain.total0.etopo1.median_pct} & \metric{rev3.marine.area_bias.marine.total0.etopo1.median_pct} & \metric{rev3.marine.area_bias.marine.marine.etopo1.median_pct} \\
SRTM15+ & \metric{h1.shape_offset.srtm15p_m} & \metric{rev3.marine.srtm15p.c_marine_m} & \metric{rev3.marine.srtm15p.shape.n_all7} & \metric{rev3.marine.srtm15p.marine.n_all7} & \metric{rev3.marine.srtm15p.shape.n_checked4} & \metric{rev3.marine.srtm15p.marine.n_checked4} & \metric{rev3.marine.area_bias.domain.total0.srtm15p.median_pct} & \metric{rev3.marine.area_bias.marine.total0.srtm15p.median_pct} & \metric{rev3.marine.area_bias.marine.marine.srtm15p.median_pct} \\
\hline
\end{tabular}
\end{table}

\emph{Area and near-term emergence.} As distributed, the grids flood \metric{rev3.marine.area_bias.marine.total0.median_min_pct}--\metric{rev3.marine.area_bias.marine.total0.median_max_pct}\% too much sea on the marine domain, against \metric{rev3.marine.area_bias.domain.total0.median_min_pct}--\metric{rev3.marine.area_bias.domain.total0.median_max_pct}\% on the whole domain. The deltaic sectors hold \metric{rev3.nearterm.S_obs.deltaic_pct}\% of the observed 2026 sea (\metric{rev3.nearterm.S_obs.deltaic_km2} of \metric{rev3.nearterm.S_obs.domain_km2}~\kmsq{}). Re-referenced to the observed 2026 shoreline on the marine domain alone (second constant \metric{rev3.nearterm.marine.c26_min_m}--\metric{rev3.nearterm.marine.c26_max_m}~m, against \metric{rev3.nearterm.registered.c26_min_m}--\metric{rev3.nearterm.registered.c26_max_m}~m), the grids place \metric{rev3.nearterm.marine.marine.E.1.min_km2}--\metric{rev3.nearterm.marine.marine.E.1.max_km2}~\kmsq{} of the observed sea in the first metre of further fall and call \metric{rev3.nearterm.marine.marine.M.min_km2}--\metric{rev3.nearterm.marine.marine.M.max_km2}~\kmsq{} of it dry. The whole-domain estimate of the article is \metric{rev3.nearterm.registered.domain.E.1.min_km2}--\metric{rev3.nearterm.registered.domain.E.1.max_km2}~\kmsq{}, of which \metric{rev3.nearterm.registered.deltaic_share_of_E1.min_pct}--\metric{rev3.nearterm.registered.deltaic_share_of_E1.max_pct}\% lies in the deltaic sectors.

\emph{Exposure counts by region.} Split with the article's constants, the first-metre counts fall mostly outside the deltaic sectors, except for protected areas. Of \metric{rev3.exposure.osm_ports.domain.min}--\metric{rev3.exposure.osm_ports.domain.max} ports or landings cut off from the sea, \metric{rev3.exposure.osm_ports.deltaic.min}--\metric{rev3.exposure.osm_ports.deltaic.max} lie in the deltaic sectors; of \metric{rev3.exposure.osm_channels.domain.min}--\metric{rev3.exposure.osm_channels.domain.max}~km of mapped channels, \metric{rev3.exposure.osm_channels.deltaic.min}--\metric{rev3.exposure.osm_channels.deltaic.max}~km; of \metric{rev3.exposure.osm_settlements.domain.min}--\metric{rev3.exposure.osm_settlements.domain.max} settlements, none. Protected areas (without the polygon drawn far larger than its reserve) hold \metric{rev3.exposure.wdpa_protected_areas_noksrz.domain.min}--\metric{rev3.exposure.wdpa_protected_areas_noksrz.domain.max}~\kmsq{} of the first metre, of which \metric{rev3.exposure.wdpa_protected_areas_noksrz.deltaic.min}--\metric{rev3.exposure.wdpa_protected_areas_noksrz.deltaic.max}~\kmsq{} lie in the deltaic sectors and \metric{rev3.exposure.wdpa_protected_areas_noksrz.marine.min}--\metric{rev3.exposure.wdpa_protected_areas_noksrz.marine.max}~\kmsq{} in the marine domain. As shares of each grid's first-metre band on the whole domain, protected areas make up \metric{rev2.n.exp.wdpa_protected_areas_noksrz.observed_reref.1.min_share}--\metric{rev2.n.exp.wdpa_protected_areas_noksrz.observed_reref.1.max_share} without that polygon and \metric{rev2.n.exp.wdpa_protected_areas.observed_reref.1.min_share}--\metric{rev2.n.exp.wdpa_protected_areas.observed_reref.1.max_share} with it. These counts are conditional on each grid and are not predictions (main article, Section~\ref{M-sec:res_nearterm}).

\emph{What the maps show.} The band that was water in \metric{params.highstand_label} and not in \metric{params.lowstand_label} covers \metric{rev3.band.domain_km2}~\kmsq{} of the domain, \metric{rev3.band.deltaic_km2}~\kmsq{} of it in the deltaic sectors. In those two sectors the maps of the article (Figs.~\ref{M-fig:rse_study_area}, \ref{M-fig:rse_emergence} and \ref{M-fig:rse_dem_accuracy}) now colour the band as emerged seabed only inside the Caspian's shoreline of February 2000. That shoreline is the Caspian polygon of HydroLAKES, version~1.0 \citep{messager2016}, whose source is the SRTM water body data and whose shorelines are those of February 2000 \citep{hydrolakes_doc}; a cell counts as sea if its centre lies inside the polygon. Landward of it, \metric{rev3.band.not_seabed_km2}~\kmsq{} (\metric{rev3.band.not_seabed_pct}\% of the domain's band; \metric{rev3.band.volga.landward_pct}\% of the Volga sector's band and \metric{rev3.band.ural.landward_pct}\% of the Ural sector's) is drawn in grey and not interpreted. Part of it may be seabed that emerged between \metric{params.highstand_label} and 2000; the data at hand cannot separate that part from the delta plain and the floodplain. This is a display rule and changes no analysis. The amendment named the raster form of that outline: the class ``saline lake'' of the Global Lakes and Wetlands Database, version~2, which was rasterised from version~1.1 of the HydroLAKES polygons and gives each \metric{rev3.band.raster.cell_arcsec} arc-second cell the wetland class with the largest share of the cell \citep{lehner2025}. That raster is too coarse for the delta fronts (the polygon has \metric{rev3.band.outline.n_vertices} vertices in the study window); with it the grey area is \metric{rev3.band.raster.not_seabed_km2}~\kmsq{} (\metric{rev3.band.raster.volga.landward_pct}\% and \metric{rev3.band.raster.ural.landward_pct}\% of the two sectors' bands), so the polygon itself is used.

\emph{The river's signature in existing data.} Two earlier results agree with the reviewer's point. The SWOT water surface of the Volga delta sector's open-sea blocks stands \metric{rev3.band.swot.volga_rel_m}~m above the surface of all sectors (\metric{rev3.band.swot.volga_n_passes} passes), while every other sector lies within \metric{rev3.band.swot.others_abs_max_m}~m of it; by calendar month the difference is \metric{rev3.band.swot.volga_month_min_m}--\metric{rev3.band.swot.volga_month_max_m}~m, from only \metric{rev3.band.swot.volga_passes_month_min}--\metric{rev3.band.swot.volga_passes_month_max} passes a month, too few to resolve the flood season. And in the late fall \metric{rev3.band.static_e3.volga_pct}\% of the sector's observed shoreline cells do not move with the level (the static channel-bank rule), against \metric{rev3.band.static_e3.others_min_pct}--\metric{rev3.band.static_e3.others_max_pct}\% in the other sectors.

\subsection*{Sector properties and a view without sector boundaries}
The seven sectors differ in area, gradient, vegetation, floor and lidar support (Supplementary Table~\ref{tab:a20_sectors}). The article therefore reports each sector's verdict on its own floor and pools no failure rate over sectors of unequal quality. To see whether the sector boundaries shape the result, we repeated the comparison on \metric{params.block_size_km}~km blocks, which ignore them. For each block and grid, $D_b$ is the median shoreline displacement pooled over all levels, and $Q_b$ the \metric{params.err_pct_hi}th percentile of the floor distances pooled the same way; a block is evaluable with at least \metric{rev3.blocks.min_n} distances in each and exceeds if $D_b > Q_b$. Of \metric{rev3.blocks.all.four.n_blocks_max} evaluable blocks, \metric{rev3.blocks.all.four.share_min_pct}--\metric{rev3.blocks.all.four.share_max_pct}\% exceed for the four failing grids and \metric{rev3.blocks.all.etopo2022.share_pct}\% for ETOPO~2022 (Supplementary Table~\ref{tab:a20_blocks}). The pattern persists in the \metric{rev3.blocks.interior.gebco2026.n_blocks} interior blocks, which touch no sector boundary (\metric{rev3.blocks.interior.four.share_min_pct}--\metric{rev3.blocks.interior.four.share_max_pct}\% exceed, against \metric{rev3.blocks.boundary.four.share_min_pct}--\metric{rev3.blocks.boundary.four.share_max_pct}\% of the \metric{rev3.blocks.boundary.gebco2026.n_blocks} blocks next to one), so the boundaries do not produce it. In the checked sectors \metric{rev3.blocks.checked.four.share_min_pct}--\metric{rev3.blocks.checked.four.share_max_pct}\% of the blocks exceed.

\begin{table}[!htbp]
\centering
\footnotesize
\setlength{\tabcolsep}{2.5pt}
\caption{\textbf{Properties of the seven sectors} (Amendment 20 item 2a). Band: the band that was water in \metric{params.highstand_label} and not in \metric{params.lowstand_label} (\metric{params.zone_scale_m}~m extents). $Q$: median over levels of the noise floor. Levels: levels with both distances (GEBCO\_2026). Gradient: median seabed gradient of the band, in metres per kilometre. Veg.: share of the band's cells dominated by emergent vegetation. Blocks: blocks of \metric{h2.block_size_km}~km with ATL06 points on bare emerged seabed, all and test. ATL06 points: their number and the effective sample size with blocks as clusters. Bias: the all-block mean error of the waterline model against ATL06. Checked sectors first.}
\label{tab:a20_sectors}
\begin{tabular}{lrrrrrrrrrrr}
\hline
 & Area & Band & Gradient & Veg. & $Q$ & & \multicolumn{2}{c}{Blocks} & \multicolumn{2}{c}{ATL06 points} & Bias \\
Sector & (\kmsq) & (\kmsq) & (m/km) & (\%) & (km) & Levels & all & test & all & effective & (m) \\
\hline
Kizlyar--Dagestan coast & \metric{study_area.sector_km2.west} & \metric{rev3.band.west.band_km2} & \metric{rev3.sector.west.gradient_m_per_km} & \metric{rev3.sector.west.vegetated_pct} & \metric{h1.noise_p95.west_km} & \metric{h1.n_evaluable.shape.gebco2026.west} & \metric{rev3.support.atl06.west.n_blocks} & \metric{rev3.support.atl06.west.n_blocks_test} & \metric{rev3.support.atl06.west.n_points} & \metric{rev3.support.atl06.west.n_eff_block} & \metric{rev2.d.obs.atl06.west.mean_m} \\
Volga--Ural interfluve coast & \metric{study_area.sector_km2.interfluve} & \metric{rev3.band.interfluve.band_km2} & \metric{rev3.sector.interfluve.gradient_m_per_km} & \metric{rev3.sector.interfluve.vegetated_pct} & \metric{h1.noise_p95.interfluve_km} & \metric{h1.n_evaluable.shape.gebco2026.interfluve} & \metric{rev3.support.atl06.interfluve.n_blocks} & \metric{rev3.support.atl06.interfluve.n_blocks_test} & \metric{rev3.support.atl06.interfluve.n_points} & \metric{rev3.support.atl06.interfluve.n_eff_block} & \metric{rev2.d.obs.atl06.interfluve.mean_m} \\
Northeastern coast & \metric{study_area.sector_km2.northeast} & \metric{rev3.band.northeast.band_km2} & \metric{rev3.sector.northeast.gradient_m_per_km} & \metric{rev3.sector.northeast.vegetated_pct} & \metric{h1.noise_p95.northeast_km} & \metric{h1.n_evaluable.shape.gebco2026.northeast} & \metric{rev3.support.atl06.northeast.n_blocks} & \metric{rev3.support.atl06.northeast.n_blocks_test} & \metric{rev3.support.atl06.northeast.n_points} & \metric{rev3.support.atl06.northeast.n_eff_block} & \metric{rev2.d.obs.atl06.northeast.mean_m} \\
Open shelf & \metric{study_area.sector_km2.open_shelf} & \metric{rev3.band.openshelf.band_km2} & \metric{rev3.sector.openshelf.gradient_m_per_km} & \metric{rev3.sector.openshelf.vegetated_pct} & \metric{h1.noise_p95.open_shelf_km} & \metric{h1.n_evaluable.shape.gebco2026.open_shelf} & \metric{rev3.support.atl06.openshelf.n_blocks} & \metric{rev3.support.atl06.openshelf.n_blocks_test} & \metric{rev3.support.atl06.openshelf.n_points} & \metric{rev3.support.atl06.openshelf.n_eff_block} & \metric{rev2.d.obs.atl06.openshelf.mean_m} \\
Volga delta (deltaic) & \metric{study_area.sector_km2.volga} & \metric{rev3.band.volga.band_km2} & \metric{rev3.sector.volga.gradient_m_per_km} & \metric{rev3.sector.volga.vegetated_pct} & \metric{h1.noise_p95.volga_km} & \metric{h1.n_evaluable.shape.gebco2026.volga} & \metric{rev3.support.atl06.volga.n_blocks} & \metric{rev3.support.atl06.volga.n_blocks_test} & \metric{rev3.support.atl06.volga.n_points} & \metric{rev3.support.atl06.volga.n_eff_block} & \metric{rev2.d.obs.atl06.volga.mean_m} \\
Ural delta (deltaic) & \metric{study_area.sector_km2.ural} & \metric{rev3.band.ural.band_km2} & \metric{rev3.sector.ural.gradient_m_per_km} & \metric{rev3.sector.ural.vegetated_pct} & \metric{h1.noise_p95.ural_km} & \metric{h1.n_evaluable.shape.gebco2026.ural} & \metric{rev3.support.atl06.ural.n_blocks} & \metric{rev3.support.atl06.ural.n_blocks_test} & \metric{rev3.support.atl06.ural.n_points} & \metric{rev3.support.atl06.ural.n_eff_block} & \metric{rev2.d.obs.atl06.ural.mean_m} \\
Komsomolets Bay & \metric{study_area.sector_km2.komsomolets} & \metric{rev3.band.komsomolets.band_km2} & \metric{rev3.sector.komsomolets.gradient_m_per_km} & \metric{rev3.sector.komsomolets.vegetated_pct} & \metric{h1.noise_p95.komsomolets_km} & \metric{h1.n_evaluable.shape.gebco2026.komsomolets} & \metric{rev3.support.atl06.komsomolets.n_blocks} & \metric{rev3.support.atl06.komsomolets.n_blocks_test} & \metric{rev3.support.atl06.komsomolets.n_points} & \metric{rev3.support.atl06.komsomolets.n_eff_block} & \metric{rev2.d.obs.atl06.komsomolets.mean_m} \\
\hline
\end{tabular}
\end{table}

\begin{table}[!htbp]
\centering
\small
\setlength{\tabcolsep}{5pt}
\caption{\textbf{Blocks whose pooled displacement exceeds their pooled floor} (Amendment 20 item 2b), as a percentage of the evaluable \metric{params.block_size_km}~km blocks; offset-corrected. Boundary blocks have a neighbouring block in another sector.}
\label{tab:a20_blocks}
\begin{tabular}{lrccccc}
\hline
Blocks & $n$ & GEBCO\_2026 & GEBCO\_2022 & ETOPO~2022 & ETOPO1 & SRTM15+ \\
\hline
All & \metric{rev3.blocks.all.gebco2026.n_blocks} & \metric{rev3.blocks.all.gebco2026.share_pct} & \metric{rev3.blocks.all.gebco2022.share_pct} & \metric{rev3.blocks.all.etopo2022.share_pct} & \metric{rev3.blocks.all.etopo1.share_pct} & \metric{rev3.blocks.all.srtm15p.share_pct} \\
Next to a sector boundary & \metric{rev3.blocks.boundary.gebco2026.n_blocks} & \metric{rev3.blocks.boundary.gebco2026.share_pct} & \metric{rev3.blocks.boundary.gebco2022.share_pct} & \metric{rev3.blocks.boundary.etopo2022.share_pct} & \metric{rev3.blocks.boundary.etopo1.share_pct} & \metric{rev3.blocks.boundary.srtm15p.share_pct} \\
Interior & \metric{rev3.blocks.interior.gebco2026.n_blocks} & \metric{rev3.blocks.interior.gebco2026.share_pct} & \metric{rev3.blocks.interior.gebco2022.share_pct} & \metric{rev3.blocks.interior.etopo2022.share_pct} & \metric{rev3.blocks.interior.etopo1.share_pct} & \metric{rev3.blocks.interior.srtm15p.share_pct} \\
Checked sectors & \metric{rev3.blocks.checked.gebco2026.n_blocks} & \metric{rev3.blocks.checked.gebco2026.share_pct} & \metric{rev3.blocks.checked.gebco2022.share_pct} & \metric{rev3.blocks.checked.etopo2022.share_pct} & \metric{rev3.blocks.checked.etopo1.share_pct} & \metric{rev3.blocks.checked.srtm15p.share_pct} \\
\quad Kizlyar--Dagestan coast & \metric{rev3.blocks.sector.west.gebco2026.n_blocks} & \metric{rev3.blocks.sector.west.gebco2026.share_pct} & \metric{rev3.blocks.sector.west.gebco2022.share_pct} & \metric{rev3.blocks.sector.west.etopo2022.share_pct} & \metric{rev3.blocks.sector.west.etopo1.share_pct} & \metric{rev3.blocks.sector.west.srtm15p.share_pct} \\
\quad Interfluve coast & \metric{rev3.blocks.sector.interfluve.gebco2026.n_blocks} & \metric{rev3.blocks.sector.interfluve.gebco2026.share_pct} & \metric{rev3.blocks.sector.interfluve.gebco2022.share_pct} & \metric{rev3.blocks.sector.interfluve.etopo2022.share_pct} & \metric{rev3.blocks.sector.interfluve.etopo1.share_pct} & \metric{rev3.blocks.sector.interfluve.srtm15p.share_pct} \\
\quad Northeastern coast & \metric{rev3.blocks.sector.northeast.gebco2026.n_blocks} & \metric{rev3.blocks.sector.northeast.gebco2026.share_pct} & \metric{rev3.blocks.sector.northeast.gebco2022.share_pct} & \metric{rev3.blocks.sector.northeast.etopo2022.share_pct} & \metric{rev3.blocks.sector.northeast.etopo1.share_pct} & \metric{rev3.blocks.sector.northeast.srtm15p.share_pct} \\
\quad Open shelf & \metric{rev3.blocks.sector.openshelf.gebco2026.n_blocks} & \metric{rev3.blocks.sector.openshelf.gebco2026.share_pct} & \metric{rev3.blocks.sector.openshelf.gebco2022.share_pct} & \metric{rev3.blocks.sector.openshelf.etopo2022.share_pct} & \metric{rev3.blocks.sector.openshelf.etopo1.share_pct} & \metric{rev3.blocks.sector.openshelf.srtm15p.share_pct} \\
Volga delta & \metric{rev3.blocks.sector.volga.gebco2026.n_blocks} & \metric{rev3.blocks.sector.volga.gebco2026.share_pct} & \metric{rev3.blocks.sector.volga.gebco2022.share_pct} & \metric{rev3.blocks.sector.volga.etopo2022.share_pct} & \metric{rev3.blocks.sector.volga.etopo1.share_pct} & \metric{rev3.blocks.sector.volga.srtm15p.share_pct} \\
Ural delta & \metric{rev3.blocks.sector.ural.gebco2026.n_blocks} & \metric{rev3.blocks.sector.ural.gebco2026.share_pct} & \metric{rev3.blocks.sector.ural.gebco2022.share_pct} & \metric{rev3.blocks.sector.ural.etopo2022.share_pct} & \metric{rev3.blocks.sector.ural.etopo1.share_pct} & \metric{rev3.blocks.sector.ural.srtm15p.share_pct} \\
Komsomolets Bay & \metric{rev3.blocks.sector.komsomolets.gebco2026.n_blocks} & \metric{rev3.blocks.sector.komsomolets.gebco2026.share_pct} & \metric{rev3.blocks.sector.komsomolets.gebco2022.share_pct} & \metric{rev3.blocks.sector.komsomolets.etopo2022.share_pct} & \metric{rev3.blocks.sector.komsomolets.etopo1.share_pct} & \metric{rev3.blocks.sector.komsomolets.srtm15p.share_pct} \\
\hline
\end{tabular}
\end{table}

\subsection*{The vertical reference that remains after the datum corrections}
Two parts must be kept apart. The first is a constant offset of the working datum. DAHITI's reference sits \metric{rev.w.v.cref.c_m_min}--\metric{rev.w.v.cref.c_m_max}~m from the ICESat-2 heights reduced with the same geoid, and the documented terms explain little of it (\metric{rev.w.v.cref.resid_documented_m_min}--\metric{rev.w.v.cref.resid_documented_m_max}~m remain); SWOT water surfaces lie \metric{rev.w.dep.common_m}~m from the DAHITI level. This constant is therefore not known to better than several decimetres. It affects only the grids read as distributed. Every verdict of the article uses grids with a fitted offset and lidar heights with a fitted constant, which absorb any constant, so no verdict depends on it. The second part is how much the reference varies from sector to sector, which no constant absorbs. It is small in the checked sectors. Their all-block mean biases of the waterline model against ICESat-2 span \metric{rev3.datum.atl06.mean_bias_range_m}~m (ATL06) and \metric{rev3.datum.atl08.mean_bias_range_m}~m (ATL08), and the largest is \metric{rev3.datum.atl06.largest_abs_mean_bias_m}~m in absolute value; their SWOT water surfaces lie within \metric{rev3.datum.swot.largest_abs_checked_m}~m of the common surface. Each sector's departure from the common reference is thus smaller than the standard deviation of the systematic level-error budget (\metric{rev.w.r.budget.sigma_median_min}--\metric{rev.w.r.budget.sigma_median_max}~m by sector). The Monte Carlo over that budget draws an offset of this size independently for each sector and epoch, and in every draw the four grids failed the Kizlyar--Dagestan, interfluve and northeastern coasts (main article, Section~\ref{M-sec:res_sep}). A reference that differs between sectors by the amounts found here therefore does not change those verdicts.

\subsection*{A lidar check of the shoreline's level, by surface type}
No manually interpreted shoreline exists for this record, and one drawn from the same Landsat images would share their limits on mud and under reeds. We used an independent measurement instead. The height of the waterline model at a point is the level at which the Landsat shoreline passed that point. Its difference from the ICESat-2 height of the same point, measured after the seabed had emerged, is therefore the error of that shoreline's level, seen by another sensor; divided by the shelf gradient it is an error of shoreline position. We computed it on all blocks, for ATL06 and ATL08 separately, by surface type and by height band (Supplementary Table~\ref{tab:a20_lidar}).

On bare seabed (emergent-vegetation share of the cell below \metric{rev3.lidar.surface_edge_lo}) of the four checked sectors, the median error is \metric{rev3.lidar.atl06.checked.bare.median_m}~m (\metric{params.ci_pct}\% interval \metric{rev3.lidar.atl06.checked.bare.median_lo_m} to \metric{rev3.lidar.atl06.checked.bare.median_hi_m}~m) and the NMAD \metric{rev3.lidar.atl06.checked.bare.nmad_m}~m, which correspond to \metric{rev3.lidar.atl06.checked.bare.median_h_abs_km} and \metric{rev3.lidar.atl06.checked.bare.nmad_h_km}~km of shoreline position. Where vegetation covers \metric{rev3.lidar.surface_edge_lo} to \metric{rev3.lidar.surface_edge_hi} of the cell the median is \metric{rev3.lidar.atl06.checked.mixed.median_m}~m. In vegetated cells (share of \metric{rev3.lidar.surface_edge_hi} or more) the model stands \metric{rev3.lidar.atl06.checked.vegetated.median_m}~m too high against ATL06 (\metric{rev3.lidar.atl06.checked.vegetated.median_lo_m} to \metric{rev3.lidar.atl06.checked.vegetated.median_hi_m}~m) and \metric{rev3.lidar.atl08.checked.vegetated.median_m}~m against ATL08: the Landsat shoreline lies seaward of the lidar ground by about \metric{rev3.lidar.atl06.checked.vegetated.median_h_abs_km}--\metric{rev3.lidar.atl08.checked.vegetated.median_h_abs_km}~km, as it must where reeds stand in shallow water and the image shows only their outer edge. Part of that difference may come from the lidar, whose ground heights under dense reeds are themselves less certain. The lowest height band (\metric{rev3.lidar.lowest_level_m} to \metric{rev3.lidar.lowest_band_top_m}~m), the seabed that emerged last and borders the present shallow water, differs by coast: the median is \metric{rev3.lidar.atl06.west.lowband.median_m}~m on the Kizlyar--Dagestan coast, \metric{rev3.lidar.atl06.interfluve.lowband.median_m}~m on the interfluve coast, \metric{rev3.lidar.atl06.northeast.lowband.median_m}~m on the northeastern coast and \metric{rev3.lidar.atl06.openshelf.lowband.median_m}~m on the open shelf. On the interfluve and northeastern coasts the lowest shorelines thus lie landward of the lidar contour, by up to \metric{rev3.lidar.both.northeast.lowband.median_h_abs_km.max}~km: recently exposed flats that surges still wet are mapped as water.

\emph{What these errors mean for the verdicts.} On the three checked coasts the largest position error of any stratum, against either product, is \metric{rev3.lidar.worst.coasts.abs_h_km.min}--\metric{rev3.lidar.worst.coasts.abs_h_km.max}~km (\metric{rev3.lidar.worst.west.stratum} cells on the Kizlyar--Dagestan and interfluve coasts, the \metric{rev3.lidar.worst.northeast.stratum} on the northeastern coast). That is \metric{rev3.lidar.worst.coasts.over_d_min.min}--\metric{rev3.lidar.worst.coasts.over_d_min.max} of the smallest sector-median displacement of the four failing grids on the same coast (main article, Section~\ref{M-sec:res_h1}). On the open shelf it is \metric{rev3.lidar.worst.openshelf.abs_h_km}~km, larger than the grids' displacements there (\metricdp{rev.s.headline.open_shelf.d_km_min}{1}--\metricdp{rev.s.headline.open_shelf.d_km_max}{1}~km), one more reason why the open-shelf verdicts carry little weight. The reed bias is larger than the level-error budget, so the budget does not cover it. Two variants of the main audit bound it instead: emergent vegetation counted as dry puts the shoreline on the seaward edge of the reed belt, and counted as water on its landward edge. Under both, each of the four grids fails all \metric{rev.s.checked.n_sectors} checked sectors (main article, Section~\ref{M-sec:res_sep}).

\begin{table}[!htbp]
\centering
\footnotesize
\setlength{\tabcolsep}{1.8pt}
\caption{\textbf{Waterline model minus ICESat-2 by surface type and for the lowest height band} (Amendment 20 item 4), over all blocks. Bare, mixed and vegetated: emergent-vegetation share of the cell below \metric{rev3.lidar.surface_edge_lo}, from \metric{rev3.lidar.surface_edge_lo} to below \metric{rev3.lidar.surface_edge_hi}, and \metric{rev3.lidar.surface_edge_hi} or more. Lowest: the lowest height band, model heights of \metric{rev3.lidar.lowest_level_m} to \metric{rev3.lidar.lowest_band_top_m}~m, on bare or mixed cells. Median with its block-bootstrap \metric{params.ci_pct}\% interval (\metric{rev3.lidar.n_boot} resamples) and NMAD, in metres. Position: the median error divided by the sector's median gradient, in kilometres (negative: the Landsat shoreline lies landward of the lidar contour). Points and blocks are those of ATL06.}
\label{tab:a20_lidar}
\begin{tabular}{llrrcrrcr}
\hline
 & & \multicolumn{5}{c}{ATL06} & \multicolumn{2}{c}{ATL08} \\
Sectors & Stratum & Points & Blocks & Median [interval] & NMAD & Position & Median [interval] & Position \\
\hline
Four checked & bare & \metric{rev3.lidar.atl06.checked.bare.n_points} & \metric{rev3.lidar.atl06.checked.bare.n_blocks} & \metric{rev3.lidar.atl06.checked.bare.median_m} [\metric{rev3.lidar.atl06.checked.bare.median_lo_m}, \metric{rev3.lidar.atl06.checked.bare.median_hi_m}] & \metric{rev3.lidar.atl06.checked.bare.nmad_m} & \metric{rev3.lidar.atl06.checked.bare.median_h_km} & \metric{rev3.lidar.atl08.checked.bare.median_m} [\metric{rev3.lidar.atl08.checked.bare.median_lo_m}, \metric{rev3.lidar.atl08.checked.bare.median_hi_m}] & \metric{rev3.lidar.atl08.checked.bare.median_h_km} \\
 & mixed & \metric{rev3.lidar.atl06.checked.mixed.n_points} & \metric{rev3.lidar.atl06.checked.mixed.n_blocks} & \metric{rev3.lidar.atl06.checked.mixed.median_m} [\metric{rev3.lidar.atl06.checked.mixed.median_lo_m}, \metric{rev3.lidar.atl06.checked.mixed.median_hi_m}] & \metric{rev3.lidar.atl06.checked.mixed.nmad_m} & \metric{rev3.lidar.atl06.checked.mixed.median_h_km} & \metric{rev3.lidar.atl08.checked.mixed.median_m} [\metric{rev3.lidar.atl08.checked.mixed.median_lo_m}, \metric{rev3.lidar.atl08.checked.mixed.median_hi_m}] & \metric{rev3.lidar.atl08.checked.mixed.median_h_km} \\
 & vegetated & \metric{rev3.lidar.atl06.checked.vegetated.n_points} & \metric{rev3.lidar.atl06.checked.vegetated.n_blocks} & \metric{rev3.lidar.atl06.checked.vegetated.median_m} [\metric{rev3.lidar.atl06.checked.vegetated.median_lo_m}, \metric{rev3.lidar.atl06.checked.vegetated.median_hi_m}] & \metric{rev3.lidar.atl06.checked.vegetated.nmad_m} & \metric{rev3.lidar.atl06.checked.vegetated.median_h_km} & \metric{rev3.lidar.atl08.checked.vegetated.median_m} [\metric{rev3.lidar.atl08.checked.vegetated.median_lo_m}, \metric{rev3.lidar.atl08.checked.vegetated.median_hi_m}] & \metric{rev3.lidar.atl08.checked.vegetated.median_h_km} \\
 & lowest &\metric{rev3.lidar.atl06.checked.lowband.n_points} & \metric{rev3.lidar.atl06.checked.lowband.n_blocks} & \metric{rev3.lidar.atl06.checked.lowband.median_m} [\metric{rev3.lidar.atl06.checked.lowband.median_lo_m}, \metric{rev3.lidar.atl06.checked.lowband.median_hi_m}] & \metric{rev3.lidar.atl06.checked.lowband.nmad_m} & \metric{rev3.lidar.atl06.checked.lowband.median_h_km} & \metric{rev3.lidar.atl08.checked.lowband.median_m} [\metric{rev3.lidar.atl08.checked.lowband.median_lo_m}, \metric{rev3.lidar.atl08.checked.lowband.median_hi_m}] & \metric{rev3.lidar.atl08.checked.lowband.median_h_km} \\
Kizlyar--Dagestan & bare & \metric{rev3.lidar.atl06.west.bare.n_points} & \metric{rev3.lidar.atl06.west.bare.n_blocks} & \metric{rev3.lidar.atl06.west.bare.median_m} [\metric{rev3.lidar.atl06.west.bare.median_lo_m}, \metric{rev3.lidar.atl06.west.bare.median_hi_m}] & \metric{rev3.lidar.atl06.west.bare.nmad_m} & \metric{rev3.lidar.atl06.west.bare.median_h_km} & \metric{rev3.lidar.atl08.west.bare.median_m} [\metric{rev3.lidar.atl08.west.bare.median_lo_m}, \metric{rev3.lidar.atl08.west.bare.median_hi_m}] & \metric{rev3.lidar.atl08.west.bare.median_h_km} \\
 & mixed & \metric{rev3.lidar.atl06.west.mixed.n_points} & \metric{rev3.lidar.atl06.west.mixed.n_blocks} & \metric{rev3.lidar.atl06.west.mixed.median_m} [\metric{rev3.lidar.atl06.west.mixed.median_lo_m}, \metric{rev3.lidar.atl06.west.mixed.median_hi_m}] & \metric{rev3.lidar.atl06.west.mixed.nmad_m} & \metric{rev3.lidar.atl06.west.mixed.median_h_km} & \metric{rev3.lidar.atl08.west.mixed.median_m} [\metric{rev3.lidar.atl08.west.mixed.median_lo_m}, \metric{rev3.lidar.atl08.west.mixed.median_hi_m}] & \metric{rev3.lidar.atl08.west.mixed.median_h_km} \\
 & vegetated & \metric{rev3.lidar.atl06.west.vegetated.n_points} & \metric{rev3.lidar.atl06.west.vegetated.n_blocks} & \metric{rev3.lidar.atl06.west.vegetated.median_m} [\metric{rev3.lidar.atl06.west.vegetated.median_lo_m}, \metric{rev3.lidar.atl06.west.vegetated.median_hi_m}] & \metric{rev3.lidar.atl06.west.vegetated.nmad_m} & \metric{rev3.lidar.atl06.west.vegetated.median_h_km} & \metric{rev3.lidar.atl08.west.vegetated.median_m} [\metric{rev3.lidar.atl08.west.vegetated.median_lo_m}, \metric{rev3.lidar.atl08.west.vegetated.median_hi_m}] & \metric{rev3.lidar.atl08.west.vegetated.median_h_km} \\
 & lowest &\metric{rev3.lidar.atl06.west.lowband.n_points} & \metric{rev3.lidar.atl06.west.lowband.n_blocks} & \metric{rev3.lidar.atl06.west.lowband.median_m} [\metric{rev3.lidar.atl06.west.lowband.median_lo_m}, \metric{rev3.lidar.atl06.west.lowband.median_hi_m}] & \metric{rev3.lidar.atl06.west.lowband.nmad_m} & \metric{rev3.lidar.atl06.west.lowband.median_h_km} & \metric{rev3.lidar.atl08.west.lowband.median_m} [\metric{rev3.lidar.atl08.west.lowband.median_lo_m}, \metric{rev3.lidar.atl08.west.lowband.median_hi_m}] & \metric{rev3.lidar.atl08.west.lowband.median_h_km} \\
Interfluve & bare & \metric{rev3.lidar.atl06.interfluve.bare.n_points} & \metric{rev3.lidar.atl06.interfluve.bare.n_blocks} & \metric{rev3.lidar.atl06.interfluve.bare.median_m} [\metric{rev3.lidar.atl06.interfluve.bare.median_lo_m}, \metric{rev3.lidar.atl06.interfluve.bare.median_hi_m}] & \metric{rev3.lidar.atl06.interfluve.bare.nmad_m} & \metric{rev3.lidar.atl06.interfluve.bare.median_h_km} & \metric{rev3.lidar.atl08.interfluve.bare.median_m} [\metric{rev3.lidar.atl08.interfluve.bare.median_lo_m}, \metric{rev3.lidar.atl08.interfluve.bare.median_hi_m}] & \metric{rev3.lidar.atl08.interfluve.bare.median_h_km} \\
 & mixed & \metric{rev3.lidar.atl06.interfluve.mixed.n_points} & \metric{rev3.lidar.atl06.interfluve.mixed.n_blocks} & \metric{rev3.lidar.atl06.interfluve.mixed.median_m} [\metric{rev3.lidar.atl06.interfluve.mixed.median_lo_m}, \metric{rev3.lidar.atl06.interfluve.mixed.median_hi_m}] & \metric{rev3.lidar.atl06.interfluve.mixed.nmad_m} & \metric{rev3.lidar.atl06.interfluve.mixed.median_h_km} & \metric{rev3.lidar.atl08.interfluve.mixed.median_m} [\metric{rev3.lidar.atl08.interfluve.mixed.median_lo_m}, \metric{rev3.lidar.atl08.interfluve.mixed.median_hi_m}] & \metric{rev3.lidar.atl08.interfluve.mixed.median_h_km} \\
 & vegetated & \metric{rev3.lidar.atl06.interfluve.vegetated.n_points} & \metric{rev3.lidar.atl06.interfluve.vegetated.n_blocks} & \metric{rev3.lidar.atl06.interfluve.vegetated.median_m} [\metric{rev3.lidar.atl06.interfluve.vegetated.median_lo_m}, \metric{rev3.lidar.atl06.interfluve.vegetated.median_hi_m}] & \metric{rev3.lidar.atl06.interfluve.vegetated.nmad_m} & \metric{rev3.lidar.atl06.interfluve.vegetated.median_h_km} & \metric{rev3.lidar.atl08.interfluve.vegetated.median_m} [\metric{rev3.lidar.atl08.interfluve.vegetated.median_lo_m}, \metric{rev3.lidar.atl08.interfluve.vegetated.median_hi_m}] & \metric{rev3.lidar.atl08.interfluve.vegetated.median_h_km} \\
 & lowest &\metric{rev3.lidar.atl06.interfluve.lowband.n_points} & \metric{rev3.lidar.atl06.interfluve.lowband.n_blocks} & \metric{rev3.lidar.atl06.interfluve.lowband.median_m} [\metric{rev3.lidar.atl06.interfluve.lowband.median_lo_m}, \metric{rev3.lidar.atl06.interfluve.lowband.median_hi_m}] & \metric{rev3.lidar.atl06.interfluve.lowband.nmad_m} & \metric{rev3.lidar.atl06.interfluve.lowband.median_h_km} & \metric{rev3.lidar.atl08.interfluve.lowband.median_m} [\metric{rev3.lidar.atl08.interfluve.lowband.median_lo_m}, \metric{rev3.lidar.atl08.interfluve.lowband.median_hi_m}] & \metric{rev3.lidar.atl08.interfluve.lowband.median_h_km} \\
Northeastern & bare & \metric{rev3.lidar.atl06.northeast.bare.n_points} & \metric{rev3.lidar.atl06.northeast.bare.n_blocks} & \metric{rev3.lidar.atl06.northeast.bare.median_m} [\metric{rev3.lidar.atl06.northeast.bare.median_lo_m}, \metric{rev3.lidar.atl06.northeast.bare.median_hi_m}] & \metric{rev3.lidar.atl06.northeast.bare.nmad_m} & \metric{rev3.lidar.atl06.northeast.bare.median_h_km} & \metric{rev3.lidar.atl08.northeast.bare.median_m} [\metric{rev3.lidar.atl08.northeast.bare.median_lo_m}, \metric{rev3.lidar.atl08.northeast.bare.median_hi_m}] & \metric{rev3.lidar.atl08.northeast.bare.median_h_km} \\
 & mixed & \metric{rev3.lidar.atl06.northeast.mixed.n_points} & \metric{rev3.lidar.atl06.northeast.mixed.n_blocks} & \metric{rev3.lidar.atl06.northeast.mixed.median_m} [\metric{rev3.lidar.atl06.northeast.mixed.median_lo_m}, \metric{rev3.lidar.atl06.northeast.mixed.median_hi_m}] & \metric{rev3.lidar.atl06.northeast.mixed.nmad_m} & \metric{rev3.lidar.atl06.northeast.mixed.median_h_km} & \metric{rev3.lidar.atl08.northeast.mixed.median_m} [\metric{rev3.lidar.atl08.northeast.mixed.median_lo_m}, \metric{rev3.lidar.atl08.northeast.mixed.median_hi_m}] & \metric{rev3.lidar.atl08.northeast.mixed.median_h_km} \\
 & vegetated & \metric{rev3.lidar.atl06.northeast.vegetated.n_points} & \metric{rev3.lidar.atl06.northeast.vegetated.n_blocks} & \metric{rev3.lidar.atl06.northeast.vegetated.median_m} [\metric{rev3.lidar.atl06.northeast.vegetated.median_lo_m}, \metric{rev3.lidar.atl06.northeast.vegetated.median_hi_m}] & \metric{rev3.lidar.atl06.northeast.vegetated.nmad_m} & \metric{rev3.lidar.atl06.northeast.vegetated.median_h_km} & \metric{rev3.lidar.atl08.northeast.vegetated.median_m} [\metric{rev3.lidar.atl08.northeast.vegetated.median_lo_m}, \metric{rev3.lidar.atl08.northeast.vegetated.median_hi_m}] & \metric{rev3.lidar.atl08.northeast.vegetated.median_h_km} \\
 & lowest &\metric{rev3.lidar.atl06.northeast.lowband.n_points} & \metric{rev3.lidar.atl06.northeast.lowband.n_blocks} & \metric{rev3.lidar.atl06.northeast.lowband.median_m} [\metric{rev3.lidar.atl06.northeast.lowband.median_lo_m}, \metric{rev3.lidar.atl06.northeast.lowband.median_hi_m}] & \metric{rev3.lidar.atl06.northeast.lowband.nmad_m} & \metric{rev3.lidar.atl06.northeast.lowband.median_h_km} & \metric{rev3.lidar.atl08.northeast.lowband.median_m} [\metric{rev3.lidar.atl08.northeast.lowband.median_lo_m}, \metric{rev3.lidar.atl08.northeast.lowband.median_hi_m}] & \metric{rev3.lidar.atl08.northeast.lowband.median_h_km} \\
Open shelf & bare & \metric{rev3.lidar.atl06.openshelf.bare.n_points} & \metric{rev3.lidar.atl06.openshelf.bare.n_blocks} & \metric{rev3.lidar.atl06.openshelf.bare.median_m} [\metric{rev3.lidar.atl06.openshelf.bare.median_lo_m}, \metric{rev3.lidar.atl06.openshelf.bare.median_hi_m}] & \metric{rev3.lidar.atl06.openshelf.bare.nmad_m} & \metric{rev3.lidar.atl06.openshelf.bare.median_h_km} & \metric{rev3.lidar.atl08.openshelf.bare.median_m} [\metric{rev3.lidar.atl08.openshelf.bare.median_lo_m}, \metric{rev3.lidar.atl08.openshelf.bare.median_hi_m}] & \metric{rev3.lidar.atl08.openshelf.bare.median_h_km} \\
 & mixed & \metric{rev3.lidar.atl06.openshelf.mixed.n_points} & \metric{rev3.lidar.atl06.openshelf.mixed.n_blocks} & \metric{rev3.lidar.atl06.openshelf.mixed.median_m} [\metric{rev3.lidar.atl06.openshelf.mixed.median_lo_m}, \metric{rev3.lidar.atl06.openshelf.mixed.median_hi_m}] & \metric{rev3.lidar.atl06.openshelf.mixed.nmad_m} & \metric{rev3.lidar.atl06.openshelf.mixed.median_h_km} & \metric{rev3.lidar.atl08.openshelf.mixed.median_m} [\metric{rev3.lidar.atl08.openshelf.mixed.median_lo_m}, \metric{rev3.lidar.atl08.openshelf.mixed.median_hi_m}] & \metric{rev3.lidar.atl08.openshelf.mixed.median_h_km} \\
 & vegetated & \metric{rev3.lidar.atl06.openshelf.vegetated.n_points} & \metric{rev3.lidar.atl06.openshelf.vegetated.n_blocks} & \metric{rev3.lidar.atl06.openshelf.vegetated.median_m} [\metric{rev3.lidar.atl06.openshelf.vegetated.median_lo_m}, \metric{rev3.lidar.atl06.openshelf.vegetated.median_hi_m}] & \metric{rev3.lidar.atl06.openshelf.vegetated.nmad_m} & \metric{rev3.lidar.atl06.openshelf.vegetated.median_h_km} & \metric{rev3.lidar.atl08.openshelf.vegetated.median_m} [\metric{rev3.lidar.atl08.openshelf.vegetated.median_lo_m}, \metric{rev3.lidar.atl08.openshelf.vegetated.median_hi_m}] & \metric{rev3.lidar.atl08.openshelf.vegetated.median_h_km} \\
 & lowest &\metric{rev3.lidar.atl06.openshelf.lowband.n_points} & \metric{rev3.lidar.atl06.openshelf.lowband.n_blocks} & \metric{rev3.lidar.atl06.openshelf.lowband.median_m} [\metric{rev3.lidar.atl06.openshelf.lowband.median_lo_m}, \metric{rev3.lidar.atl06.openshelf.lowband.median_hi_m}] & \metric{rev3.lidar.atl06.openshelf.lowband.nmad_m} & \metric{rev3.lidar.atl06.openshelf.lowband.median_h_km} & \metric{rev3.lidar.atl08.openshelf.lowband.median_m} [\metric{rev3.lidar.atl08.openshelf.lowband.median_lo_m}, \metric{rev3.lidar.atl08.openshelf.lowband.median_hi_m}] & \metric{rev3.lidar.atl08.openshelf.lowband.median_h_km} \\
\hline
\end{tabular}
\end{table}

\subsection*{Repeatability, accuracy and the independence of the accuracy assessment}
The noise floor $Q$ measures repeatability: how closely two interleaved halves of the same imagery place the same shoreline. The halves share the level series, the threshold method and the vegetation mask, so $Q$ says nothing about errors common to both. Absolute accuracy comes from the lidar comparison above and from the all-block biases (Supplementary Table~\ref{tab:a19_obs}), and the errors common to all observations are bounded by the level-error budget.

Spatial leakage would arise if lidar heights near a test point had shaped the model that is tested there. The waterline model is built from Landsat waterlines and altimetric levels alone; no lidar height enters it. Its one tuned quantity, the threshold of the static channel-bank rule, is set from the water frequencies of training blocks and uses no lidar. Lidar heights fit only the datum constant of each lidar product and one constant per grid, on training blocks, and the scale and quantiles of the conformal intervals, on training and calibration blocks. Test-block points fit nothing. In the hindcast the only fitted quantity is each grid's offset $c$, fitted on the training blocks, and the out-of-sample test of that offset is scored on the calibration and test blocks. One use of the test blocks is not blind: the checked-sector criterion was set after the per-sector test-block errors had been seen (main article, Section~\ref{M-sec:dem}). The biases over all blocks, most of which played no part in that choice, are the safeguard (Supplementary Table~\ref{tab:a19_obs}). As a direct test we recomputed the test-block statistics on the points that lie at least \metric{rev3.indep.buffer_km}~km from every training or calibration block (\metric{rev3.indep.atl08.buf10.share_pct}\% of the ATL08 test points; Supplementary Table~\ref{tab:a20_indep}). The model's NMAD is \metric{rev3.indep.atl08.buf10.dem.nmad_m}~m there against \metric{rev3.indep.atl08.all.dem.nmad_m}~m on all test points, and the grids' \metric{rev3.indep.atl08.buf10.grids.nmad_min_m}--\metric{rev3.indep.atl08.buf10.grids.nmad_max_m}~m against \metric{rev3.indep.atl08.all.grids.nmad_min_m}--\metric{rev3.indep.atl08.all.grids.nmad_max_m}~m. The model's mean absolute error within \metric{rev3.indep.buffer_km}~km of a training or calibration block differs from that farther away by \metric{rev3.indep.atl08.near_far.difference_m}~m (\metric{params.ci_pct}\% interval \metric{rev3.indep.atl08.near_far.lo_m} to \metric{rev3.indep.atl08.near_far.hi_m}~m). The interval includes zero, so the test detects no advantage near the training blocks; with \metric{rev3.indep.atl08.all.n_blocks} test blocks it could not detect a small one. No test point lies \metric{rev3.indep.buffer_far_km}~km or more from such a block, so the second pre-specified buffer could not be evaluated.

\begin{table}[!htbp]
\centering
\small
\caption{\textbf{Test-block accuracy at a distance from the training and calibration blocks} (Amendment 20 item 5). NMAD and RMSE in metres; grids: range over the five grids, each with its training-block offset. Nothing is refitted.}
\label{tab:a20_indep}
\begin{tabular}{llrrcccc}
\hline
 & & & & \multicolumn{2}{c}{Waterline model} & \multicolumn{2}{c}{Grids} \\
Reference & Test points & Points & Blocks & NMAD & RMSE & NMAD & RMSE \\
\hline
ATL08 & all & \metric{rev3.indep.atl08.all.n_points} & \metric{rev3.indep.atl08.all.n_blocks} & \metric{rev3.indep.atl08.all.dem.nmad_m} & \metric{rev3.indep.atl08.all.dem.rmse_m} & \metric{rev3.indep.atl08.all.grids.nmad_min_m}--\metric{rev3.indep.atl08.all.grids.nmad_max_m} & \metric{rev3.indep.atl08.all.grids.rmse_min_m}--\metric{rev3.indep.atl08.all.grids.rmse_max_m} \\
ATL08 & at least \metric{rev3.indep.buffer_km}~km away & \metric{rev3.indep.atl08.buf10.n_points} & \metric{rev3.indep.atl08.buf10.n_blocks} & \metric{rev3.indep.atl08.buf10.dem.nmad_m} & \metric{rev3.indep.atl08.buf10.dem.rmse_m} & \metric{rev3.indep.atl08.buf10.grids.nmad_min_m}--\metric{rev3.indep.atl08.buf10.grids.nmad_max_m} & \metric{rev3.indep.atl08.buf10.grids.rmse_min_m}--\metric{rev3.indep.atl08.buf10.grids.rmse_max_m} \\
ATL06 & all & \metric{rev3.indep.atl06.all.n_points} & \metric{rev3.indep.atl06.all.n_blocks} & \metric{rev3.indep.atl06.all.dem.nmad_m} & \metric{rev3.indep.atl06.all.dem.rmse_m} & \metric{rev3.indep.atl06.all.grids.nmad_min_m}--\metric{rev3.indep.atl06.all.grids.nmad_max_m} & \metric{rev3.indep.atl06.all.grids.rmse_min_m}--\metric{rev3.indep.atl06.all.grids.rmse_max_m} \\
ATL06 & at least \metric{rev3.indep.buffer_km}~km away & \metric{rev3.indep.atl06.buf10.n_points} & \metric{rev3.indep.atl06.buf10.n_blocks} & \metric{rev3.indep.atl06.buf10.dem.nmad_m} & \metric{rev3.indep.atl06.buf10.dem.rmse_m} & \metric{rev3.indep.atl06.buf10.grids.nmad_min_m}--\metric{rev3.indep.atl06.buf10.grids.nmad_max_m} & \metric{rev3.indep.atl06.buf10.grids.rmse_min_m}--\metric{rev3.indep.atl06.buf10.grids.rmse_max_m} \\
\hline
\end{tabular}
\end{table}

\subsection*{The practical size of the misses}
The decision rule compares a median displacement with the floor, so on a quiet coast it can flag a displacement that is real but small. To separate such cases we laid three further thresholds over the rule, each fixed earlier for another purpose: the displacement of the degraded-DEM control at the same level, which is what whole-metre storage and the grid's resolution alone produce; one native cell of the grid (\metric{rev3.tier.cell_15s_km}~km at 15 arc-second resolution, \metric{rev3.tier.cell_1m_km}~km at 1 arc-minute); and the distance that \metric{rev3.tier.height_m}~m of height spans at the sector's median gradient (\metric{rev3.tier.half_metre_km.coasts.min}--\metric{rev3.tier.half_metre_km.coasts.max}~km on the three checked coasts, \metric{rev3.tier.half_metre_km.openshelf}~km on the open shelf), the tolerance of the checked-sector criterion. A level is \emph{material} if the displacement exceeds the floor and all three; a grid's failure in a sector is material if that holds at more than half of the levels, and \emph{resolution-scale} if the rule fails without it. The classes were defined after the displacements were known and decide nothing.

Of the \metric{rev3.tier.four.n_cases} failures of the four grids in the checked sectors, \metric{rev3.tier.four.n_material} are material and \metric{rev3.tier.four.n_resolution} are resolution-scale; the latter are the four open-shelf cases (Supplementary Table~\ref{tab:a20_tiers}). On the Kizlyar--Dagestan and interfluve coasts the material majority holds in \metric{rev3.tier.four.westinterfluve.p_material.min} to \metric{rev3.tier.four.westinterfluve.p_material.max} of the level resamples, and the median excess of the displacement over the floor has lower interval bounds of \metric{rev3.tier.four.westinterfluve.d_minus_q_lo_km.min}--\metric{rev3.tier.four.westinterfluve.d_minus_q_lo_km.max}~km. On the northeastern coast the failures are material but less stable (\metric{rev3.tier.four.northeast.p_material.min}--\metric{rev3.tier.four.northeast.p_material.max} of the resamples).

\begin{table}[!htbp]
\centering
\footnotesize
\setlength{\tabcolsep}{3.5pt}
\caption{\textbf{Practical size of each grid's misses in the checked sectors} (Amendment 20 item 6; an overlay defined after the results, which decides nothing). Levels: above the floor / material / with both distances. $D - Q$: median over levels of the displacement less the floor, in kilometres, with its moving-block bootstrap \metric{params.ci_pct}\% interval (blocks of \metric{rev3.tier.mbb_block_levels} levels, \metric{rev3.tier.n_boot} resamples). Share: resamples with a material majority.}
\label{tab:a20_tiers}
\begin{tabular}{llllrr}
\hline
Grid & Sector & Class & Levels & $D - Q$ (km) [interval] & Share \\
\hline
GEBCO\_2026 & Kizlyar--Dagestan & \metric{rev3.tier.gebco2026.west.class} & \metric{rev3.tier.gebco2026.west.n_above_floor} / \metric{rev3.tier.gebco2026.west.n_material} / \metric{rev3.tier.gebco2026.west.n_evaluable} & \metric{rev3.tier.gebco2026.west.d_minus_q_km} [\metric{rev3.tier.gebco2026.west.d_minus_q_lo_km}, \metric{rev3.tier.gebco2026.west.d_minus_q_hi_km}] & \metric{rev3.tier.gebco2026.west.p_material} \\
 & Interfluve & \metric{rev3.tier.gebco2026.interfluve.class} & \metric{rev3.tier.gebco2026.interfluve.n_above_floor} / \metric{rev3.tier.gebco2026.interfluve.n_material} / \metric{rev3.tier.gebco2026.interfluve.n_evaluable} & \metric{rev3.tier.gebco2026.interfluve.d_minus_q_km} [\metric{rev3.tier.gebco2026.interfluve.d_minus_q_lo_km}, \metric{rev3.tier.gebco2026.interfluve.d_minus_q_hi_km}] & \metric{rev3.tier.gebco2026.interfluve.p_material} \\
 & Northeastern coast & \metric{rev3.tier.gebco2026.northeast.class} & \metric{rev3.tier.gebco2026.northeast.n_above_floor} / \metric{rev3.tier.gebco2026.northeast.n_material} / \metric{rev3.tier.gebco2026.northeast.n_evaluable} & \metric{rev3.tier.gebco2026.northeast.d_minus_q_km} [\metric{rev3.tier.gebco2026.northeast.d_minus_q_lo_km}, \metric{rev3.tier.gebco2026.northeast.d_minus_q_hi_km}] & \metric{rev3.tier.gebco2026.northeast.p_material} \\
 & Open shelf & \metric{rev3.tier.gebco2026.openshelf.class} & \metric{rev3.tier.gebco2026.openshelf.n_above_floor} / \metric{rev3.tier.gebco2026.openshelf.n_material} / \metric{rev3.tier.gebco2026.openshelf.n_evaluable} & \metric{rev3.tier.gebco2026.openshelf.d_minus_q_km} [\metric{rev3.tier.gebco2026.openshelf.d_minus_q_lo_km}, \metric{rev3.tier.gebco2026.openshelf.d_minus_q_hi_km}] & \metric{rev3.tier.gebco2026.openshelf.p_material} \\
GEBCO\_2022 & Kizlyar--Dagestan & \metric{rev3.tier.gebco2022.west.class} & \metric{rev3.tier.gebco2022.west.n_above_floor} / \metric{rev3.tier.gebco2022.west.n_material} / \metric{rev3.tier.gebco2022.west.n_evaluable} & \metric{rev3.tier.gebco2022.west.d_minus_q_km} [\metric{rev3.tier.gebco2022.west.d_minus_q_lo_km}, \metric{rev3.tier.gebco2022.west.d_minus_q_hi_km}] & \metric{rev3.tier.gebco2022.west.p_material} \\
 & Interfluve & \metric{rev3.tier.gebco2022.interfluve.class} & \metric{rev3.tier.gebco2022.interfluve.n_above_floor} / \metric{rev3.tier.gebco2022.interfluve.n_material} / \metric{rev3.tier.gebco2022.interfluve.n_evaluable} & \metric{rev3.tier.gebco2022.interfluve.d_minus_q_km} [\metric{rev3.tier.gebco2022.interfluve.d_minus_q_lo_km}, \metric{rev3.tier.gebco2022.interfluve.d_minus_q_hi_km}] & \metric{rev3.tier.gebco2022.interfluve.p_material} \\
 & Northeastern coast & \metric{rev3.tier.gebco2022.northeast.class} & \metric{rev3.tier.gebco2022.northeast.n_above_floor} / \metric{rev3.tier.gebco2022.northeast.n_material} / \metric{rev3.tier.gebco2022.northeast.n_evaluable} & \metric{rev3.tier.gebco2022.northeast.d_minus_q_km} [\metric{rev3.tier.gebco2022.northeast.d_minus_q_lo_km}, \metric{rev3.tier.gebco2022.northeast.d_minus_q_hi_km}] & \metric{rev3.tier.gebco2022.northeast.p_material} \\
 & Open shelf & \metric{rev3.tier.gebco2022.openshelf.class} & \metric{rev3.tier.gebco2022.openshelf.n_above_floor} / \metric{rev3.tier.gebco2022.openshelf.n_material} / \metric{rev3.tier.gebco2022.openshelf.n_evaluable} & \metric{rev3.tier.gebco2022.openshelf.d_minus_q_km} [\metric{rev3.tier.gebco2022.openshelf.d_minus_q_lo_km}, \metric{rev3.tier.gebco2022.openshelf.d_minus_q_hi_km}] & \metric{rev3.tier.gebco2022.openshelf.p_material} \\
ETOPO~2022 & Kizlyar--Dagestan & \metric{rev3.tier.etopo2022.west.class} & \metric{rev3.tier.etopo2022.west.n_above_floor} / \metric{rev3.tier.etopo2022.west.n_material} / \metric{rev3.tier.etopo2022.west.n_evaluable} & \metric{rev3.tier.etopo2022.west.d_minus_q_km} [\metric{rev3.tier.etopo2022.west.d_minus_q_lo_km}, \metric{rev3.tier.etopo2022.west.d_minus_q_hi_km}] & \metric{rev3.tier.etopo2022.west.p_material} \\
 & Interfluve & \metric{rev3.tier.etopo2022.interfluve.class} & \metric{rev3.tier.etopo2022.interfluve.n_above_floor} / \metric{rev3.tier.etopo2022.interfluve.n_material} / \metric{rev3.tier.etopo2022.interfluve.n_evaluable} & \metric{rev3.tier.etopo2022.interfluve.d_minus_q_km} [\metric{rev3.tier.etopo2022.interfluve.d_minus_q_lo_km}, \metric{rev3.tier.etopo2022.interfluve.d_minus_q_hi_km}] & \metric{rev3.tier.etopo2022.interfluve.p_material} \\
 & Northeastern coast & \metric{rev3.tier.etopo2022.northeast.class} & \metric{rev3.tier.etopo2022.northeast.n_above_floor} / \metric{rev3.tier.etopo2022.northeast.n_material} / \metric{rev3.tier.etopo2022.northeast.n_evaluable} & \metric{rev3.tier.etopo2022.northeast.d_minus_q_km} [\metric{rev3.tier.etopo2022.northeast.d_minus_q_lo_km}, \metric{rev3.tier.etopo2022.northeast.d_minus_q_hi_km}] & \metric{rev3.tier.etopo2022.northeast.p_material} \\
 & Open shelf & \metric{rev3.tier.etopo2022.openshelf.class} & \metric{rev3.tier.etopo2022.openshelf.n_above_floor} / \metric{rev3.tier.etopo2022.openshelf.n_material} / \metric{rev3.tier.etopo2022.openshelf.n_evaluable} & \metric{rev3.tier.etopo2022.openshelf.d_minus_q_km} [\metric{rev3.tier.etopo2022.openshelf.d_minus_q_lo_km}, \metric{rev3.tier.etopo2022.openshelf.d_minus_q_hi_km}] & \metric{rev3.tier.etopo2022.openshelf.p_material} \\
ETOPO1 & Kizlyar--Dagestan & \metric{rev3.tier.etopo1.west.class} & \metric{rev3.tier.etopo1.west.n_above_floor} / \metric{rev3.tier.etopo1.west.n_material} / \metric{rev3.tier.etopo1.west.n_evaluable} & \metric{rev3.tier.etopo1.west.d_minus_q_km} [\metric{rev3.tier.etopo1.west.d_minus_q_lo_km}, \metric{rev3.tier.etopo1.west.d_minus_q_hi_km}] & \metric{rev3.tier.etopo1.west.p_material} \\
 & Interfluve & \metric{rev3.tier.etopo1.interfluve.class} & \metric{rev3.tier.etopo1.interfluve.n_above_floor} / \metric{rev3.tier.etopo1.interfluve.n_material} / \metric{rev3.tier.etopo1.interfluve.n_evaluable} & \metric{rev3.tier.etopo1.interfluve.d_minus_q_km} [\metric{rev3.tier.etopo1.interfluve.d_minus_q_lo_km}, \metric{rev3.tier.etopo1.interfluve.d_minus_q_hi_km}] & \metric{rev3.tier.etopo1.interfluve.p_material} \\
 & Northeastern coast & \metric{rev3.tier.etopo1.northeast.class} & \metric{rev3.tier.etopo1.northeast.n_above_floor} / \metric{rev3.tier.etopo1.northeast.n_material} / \metric{rev3.tier.etopo1.northeast.n_evaluable} & \metric{rev3.tier.etopo1.northeast.d_minus_q_km} [\metric{rev3.tier.etopo1.northeast.d_minus_q_lo_km}, \metric{rev3.tier.etopo1.northeast.d_minus_q_hi_km}] & \metric{rev3.tier.etopo1.northeast.p_material} \\
 & Open shelf & \metric{rev3.tier.etopo1.openshelf.class} & \metric{rev3.tier.etopo1.openshelf.n_above_floor} / \metric{rev3.tier.etopo1.openshelf.n_material} / \metric{rev3.tier.etopo1.openshelf.n_evaluable} & \metric{rev3.tier.etopo1.openshelf.d_minus_q_km} [\metric{rev3.tier.etopo1.openshelf.d_minus_q_lo_km}, \metric{rev3.tier.etopo1.openshelf.d_minus_q_hi_km}] & \metric{rev3.tier.etopo1.openshelf.p_material} \\
SRTM15+ & Kizlyar--Dagestan & \metric{rev3.tier.srtm15p.west.class} & \metric{rev3.tier.srtm15p.west.n_above_floor} / \metric{rev3.tier.srtm15p.west.n_material} / \metric{rev3.tier.srtm15p.west.n_evaluable} & \metric{rev3.tier.srtm15p.west.d_minus_q_km} [\metric{rev3.tier.srtm15p.west.d_minus_q_lo_km}, \metric{rev3.tier.srtm15p.west.d_minus_q_hi_km}] & \metric{rev3.tier.srtm15p.west.p_material} \\
 & Interfluve & \metric{rev3.tier.srtm15p.interfluve.class} & \metric{rev3.tier.srtm15p.interfluve.n_above_floor} / \metric{rev3.tier.srtm15p.interfluve.n_material} / \metric{rev3.tier.srtm15p.interfluve.n_evaluable} & \metric{rev3.tier.srtm15p.interfluve.d_minus_q_km} [\metric{rev3.tier.srtm15p.interfluve.d_minus_q_lo_km}, \metric{rev3.tier.srtm15p.interfluve.d_minus_q_hi_km}] & \metric{rev3.tier.srtm15p.interfluve.p_material} \\
 & Northeastern coast & \metric{rev3.tier.srtm15p.northeast.class} & \metric{rev3.tier.srtm15p.northeast.n_above_floor} / \metric{rev3.tier.srtm15p.northeast.n_material} / \metric{rev3.tier.srtm15p.northeast.n_evaluable} & \metric{rev3.tier.srtm15p.northeast.d_minus_q_km} [\metric{rev3.tier.srtm15p.northeast.d_minus_q_lo_km}, \metric{rev3.tier.srtm15p.northeast.d_minus_q_hi_km}] & \metric{rev3.tier.srtm15p.northeast.p_material} \\
 & Open shelf & \metric{rev3.tier.srtm15p.openshelf.class} & \metric{rev3.tier.srtm15p.openshelf.n_above_floor} / \metric{rev3.tier.srtm15p.openshelf.n_material} / \metric{rev3.tier.srtm15p.openshelf.n_evaluable} & \metric{rev3.tier.srtm15p.openshelf.d_minus_q_km} [\metric{rev3.tier.srtm15p.openshelf.d_minus_q_lo_km}, \metric{rev3.tier.srtm15p.openshelf.d_minus_q_hi_km}] & \metric{rev3.tier.srtm15p.openshelf.p_material} \\
\hline
\end{tabular}
\end{table}

\subsection*{ICESat-2 support by sector}
The lidar points are many but strongly clustered, so their number overstates the information they carry (Supplementary Table~\ref{tab:a20_support}). Within a sector the errors of points in one \metric{h2.block_size_km}~km block are correlated; with blocks as clusters, the effective sample size of the checked sectors is \metric{rev3.support.atl06.checked.n_eff_block.min}--\metric{rev3.support.atl06.checked.n_eff_block.max} points against \metric{rev3.support.atl06.checked.n_points.min}--\metric{rev3.support.atl06.checked.n_points.max} ATL06 points. This is why every interval in the article resamples blocks, not points, and why the Kizlyar--Dagestan and interfluve coasts, with one test block each, are also judged on all of their blocks. Supplementary Fig.~\ref{fig:a20_lidar_blocks} maps the blocks, their split, their points and the model's mean error in each.

\begin{figure}[!htbp]
  \centering
  \includegraphics[width=\linewidth]{../figures/out/rse_si_fig_lidar_blocks.pdf}
  \caption{\textbf{Where the ICESat-2 support lies} (Amendment 20 item 7). The \metric{rev3.support.atl06.overall.n_blocks} blocks of \metric{h2.block_size_km}~km that hold ICESat-2 ATL06 points on bare emerged seabed, outlined by their split (training, calibration or test). Number in a block: its points, in thousands. Fill: the mean error of the waterline model in the block, model minus ICESat-2 (the colour scale is clipped; see the key). Thin lines: the sectors, with the checked sectors named in bold. Light blue: the open sea of the \metric{params.lowstand_label} low stand. Map lines delineate study areas and do not necessarily depict accepted national boundaries.}
  \label{fig:a20_lidar_blocks}
\end{figure}

\begin{table}[!htbp]
\centering
\footnotesize
\setlength{\tabcolsep}{3.5pt}
\caption{\textbf{ICESat-2 ATL06 support of each sector} (Amendment 20 item 7), on bare emerged seabed, over all blocks of \metric{h2.block_size_km}~km. Blocks: all (training / calibration / test). Tracks: track segments (granule and beam). Correlation: intra-block correlation of the model's error. Effective points: Kish's design effect with blocks, and with track segments, as clusters. The ATL08 table is in the data release.}
\label{tab:a20_support}
\begin{tabular}{lrlrrcrrr}
\hline
 & & & & & & & \multicolumn{2}{c}{Effective points} \\
Sector & Points & Blocks & Tracks & Dates & Years & Correlation & blocks & tracks \\
\hline
Kizlyar--Dagestan coast & \metric{rev3.support.atl06.west.n_points} & \metric{rev3.support.atl06.west.n_blocks} (\metric{rev3.support.atl06.west.n_blocks_train} / \metric{rev3.support.atl06.west.n_blocks_calib} / \metric{rev3.support.atl06.west.n_blocks_test}) & \metric{rev3.support.atl06.west.n_segments} & \metric{rev3.support.atl06.west.n_dates} & \metric{rev3.support.atl06.west.first_year}--\metric{rev3.support.atl06.west.last_year} & \metric{rev3.support.atl06.west.icc_block} & \metric{rev3.support.atl06.west.n_eff_block} & \metric{rev3.support.atl06.west.n_eff_segment} \\
Volga--Ural interfluve coast & \metric{rev3.support.atl06.interfluve.n_points} & \metric{rev3.support.atl06.interfluve.n_blocks} (\metric{rev3.support.atl06.interfluve.n_blocks_train} / \metric{rev3.support.atl06.interfluve.n_blocks_calib} / \metric{rev3.support.atl06.interfluve.n_blocks_test}) & \metric{rev3.support.atl06.interfluve.n_segments} & \metric{rev3.support.atl06.interfluve.n_dates} & \metric{rev3.support.atl06.interfluve.first_year}--\metric{rev3.support.atl06.interfluve.last_year} & \metric{rev3.support.atl06.interfluve.icc_block} & \metric{rev3.support.atl06.interfluve.n_eff_block} & \metric{rev3.support.atl06.interfluve.n_eff_segment} \\
Northeastern coast & \metric{rev3.support.atl06.northeast.n_points} & \metric{rev3.support.atl06.northeast.n_blocks} (\metric{rev3.support.atl06.northeast.n_blocks_train} / \metric{rev3.support.atl06.northeast.n_blocks_calib} / \metric{rev3.support.atl06.northeast.n_blocks_test}) & \metric{rev3.support.atl06.northeast.n_segments} & \metric{rev3.support.atl06.northeast.n_dates} & \metric{rev3.support.atl06.northeast.first_year}--\metric{rev3.support.atl06.northeast.last_year} & \metric{rev3.support.atl06.northeast.icc_block} & \metric{rev3.support.atl06.northeast.n_eff_block} & \metric{rev3.support.atl06.northeast.n_eff_segment} \\
Open shelf & \metric{rev3.support.atl06.openshelf.n_points} & \metric{rev3.support.atl06.openshelf.n_blocks} (\metric{rev3.support.atl06.openshelf.n_blocks_train} / \metric{rev3.support.atl06.openshelf.n_blocks_calib} / \metric{rev3.support.atl06.openshelf.n_blocks_test}) & \metric{rev3.support.atl06.openshelf.n_segments} & \metric{rev3.support.atl06.openshelf.n_dates} & \metric{rev3.support.atl06.openshelf.first_year}--\metric{rev3.support.atl06.openshelf.last_year} & \metric{rev3.support.atl06.openshelf.icc_block} & \metric{rev3.support.atl06.openshelf.n_eff_block} & \metric{rev3.support.atl06.openshelf.n_eff_segment} \\
Volga delta & \metric{rev3.support.atl06.volga.n_points} & \metric{rev3.support.atl06.volga.n_blocks} (\metric{rev3.support.atl06.volga.n_blocks_train} / \metric{rev3.support.atl06.volga.n_blocks_calib} / \metric{rev3.support.atl06.volga.n_blocks_test}) & \metric{rev3.support.atl06.volga.n_segments} & \metric{rev3.support.atl06.volga.n_dates} & \metric{rev3.support.atl06.volga.first_year}--\metric{rev3.support.atl06.volga.last_year} & \metric{rev3.support.atl06.volga.icc_block} & \metric{rev3.support.atl06.volga.n_eff_block} & \metric{rev3.support.atl06.volga.n_eff_segment} \\
Ural delta & \metric{rev3.support.atl06.ural.n_points} & \metric{rev3.support.atl06.ural.n_blocks} (\metric{rev3.support.atl06.ural.n_blocks_train} / \metric{rev3.support.atl06.ural.n_blocks_calib} / \metric{rev3.support.atl06.ural.n_blocks_test}) & \metric{rev3.support.atl06.ural.n_segments} & \metric{rev3.support.atl06.ural.n_dates} & \metric{rev3.support.atl06.ural.first_year}--\metric{rev3.support.atl06.ural.last_year} & \metric{rev3.support.atl06.ural.icc_block} & \metric{rev3.support.atl06.ural.n_eff_block} & \metric{rev3.support.atl06.ural.n_eff_segment} \\
Komsomolets Bay & \metric{rev3.support.atl06.komsomolets.n_points} & \metric{rev3.support.atl06.komsomolets.n_blocks} (\metric{rev3.support.atl06.komsomolets.n_blocks_train} / \metric{rev3.support.atl06.komsomolets.n_blocks_calib} / \metric{rev3.support.atl06.komsomolets.n_blocks_test}) & \metric{rev3.support.atl06.komsomolets.n_segments} & \metric{rev3.support.atl06.komsomolets.n_dates} & \metric{rev3.support.atl06.komsomolets.first_year}--\metric{rev3.support.atl06.komsomolets.last_year} & \metric{rev3.support.atl06.komsomolets.icc_block} & \metric{rev3.support.atl06.komsomolets.n_eff_block} & \metric{rev3.support.atl06.komsomolets.n_eff_segment} \\
All sectors & \metric{rev3.support.atl06.overall.n_points} & \metric{rev3.support.atl06.overall.n_blocks} (\metric{rev3.support.atl06.overall.n_blocks_train} / \metric{rev3.support.atl06.overall.n_blocks_calib} / \metric{rev3.support.atl06.overall.n_blocks_test}) & \metric{rev3.support.atl06.overall.n_segments} & \metric{rev3.support.atl06.overall.n_dates} & \metric{rev3.support.atl06.overall.first_year}--\metric{rev3.support.atl06.overall.last_year} & \metric{rev3.support.atl06.overall.icc_block} & \metric{rev3.support.atl06.overall.n_eff_block} & \metric{rev3.support.atl06.overall.n_eff_segment} \\
\hline
\end{tabular}
\end{table}

\subsection*{Wind setup at the acquisition times}
\emph{Why the screening removed no scene.} The rule did not find the acquisitions calm; it found its own statistic uninformative. That statistic is the median distance between the two halves' shorelines, and it stood at one \metric{params.stage2_res_m}~m grid cell without screening and at every threshold, while each stricter threshold discarded usable data (Section~\SImethods, Wind-setup screening). Most of each shoreline does not move by a cell between the halves; setup acts on the tail of the distances, which a median does not see and which the floor, a \metric{params.err_pct_hi}th percentile, does.

\emph{How large setup is.} For the \metric{rev3.setup.n_acquisitions} Landsat acquisitions of the main processing plan, the predicted local departure from the basin level exceeds \metric{rev3.setup.threshold_m}~m in absolute value in \metric{rev3.setup.coasts.share_above_0p20_pct.min}--\metric{rev3.setup.coasts.share_above_0p20_pct.max}\% of the acquisitions on the three checked coasts (Supplementary Table~\ref{tab:a20_setup}). The acquisitions are not unusually calm: the median ERA5 wind speed at the acquisition hours is \metric{rev3.setup.wind.ratio_same_hours.min}--\metric{rev3.setup.wind.ratio_same_hours.max} times that of the same hours of the day in the same months and years (\metric{rev3.setup.wind.ratio_all_hours.min}--\metric{rev3.setup.wind.ratio_all_hours.max} times that of all hours, because the morning overpass falls in the windier part of the day). With the model's calibration residual added, the departure of one acquisition has a standard deviation of \metric{rev3.setup.coasts.total_sd_m.min}--\metric{rev3.setup.coasts.total_sd_m.max}~m on these coasts, which moves the shoreline by \metric{rev3.setup.coasts.total_sd_h_km.min}--\metric{rev3.setup.coasts.total_sd_h_km.max}~km at the median gradient, \metric{rev3.setup.coasts.total_sd_over_floor.min}--\metric{rev3.setup.coasts.total_sd_over_floor.max} of the floor.

\emph{What it does to a level-binned shoreline.} A shoreline of the audit is not one acquisition. It is the contour on which the cells are water in half of a level bin's acquisitions, so a departure that varies from one acquisition to the next widens the floor, as intended, and shifts the shoreline little. What can shift it is a departure that the acquisitions of a bin share. The mean predicted departure of a bin has a median of \metric{rev3.setup.coasts.unit_mean_abs_median_m.min}--\metric{rev3.setup.coasts.unit_mean_abs_median_m.max}~m in absolute value on the three checked coasts and reaches \metric{rev3.setup.coasts.unit_mean_abs_max_m.min}--\metric{rev3.setup.coasts.unit_mean_abs_max_m.max}~m in the worst bin of each, or \metric{rev3.setup.coasts.unit_mean_abs_max_h_km.min}--\metric{rev3.setup.coasts.unit_mean_abs_max_h_km.max}~km of shoreline. The worst bin is \metric{rev3.setup.coasts.unit_mean_abs_max_over_sigma_b.min}--\metric{rev3.setup.coasts.unit_mean_abs_max_over_sigma_b.max} times the standard deviation of the coast's systematic level-error budget, which already contains the setup model's residual. The Monte Carlo over that budget shifts all levels of a sector and epoch together by offsets of this size, and in every draw the four grids failed the three coasts (main article, Section~\ref{M-sec:res_sep}). The predicted setup is a model estimate with a residual of its own, so these numbers bound the effect and do not measure it.

\emph{Earlier direct tests.} Two variants of the processing had already tested the effect of setup on the observations (Supplementary Table~\ref{tab:si_robust}, check~4). Screening at a stricter fixed threshold, on a sample of test blocks, and correcting each scene's level by its predicted setup in place of screening both left the accuracy of the waterline model unchanged (ATL08 NMAD \metricdp{robustness.check04.setup_correction.atl08.dem_nmad_m}{2}~m with the correction and \metricdp{robustness.check04.setup_correction.atl08.dem_nmad_m_primary}{2}~m without; each variant was judged ``\metric{robustness.check04.setup_correction.verdict}'' by the rule of that section). Neither variant re-ran the hindcast audit. A split of the scenes by wind regime, which would do so, needs new image processing and was not run.

\begin{table}[!htbp]
\centering
\small
\setlength{\tabcolsep}{4pt}
\caption{\textbf{Predicted wind setup at the Landsat acquisition times} (Amendment 20 item 8), by sector. SD and the \metric{params.err_pct_hi}th percentile of the absolute predicted departure from the basin level; share of acquisitions with an absolute departure above \metric{rev3.setup.threshold_m}~m; total SD with the model's calibration residual added in quadrature, in metres and as shoreline distance at the median gradient; that distance over the sector's median floor; and the largest absolute mean departure of a level bin. Komsomolets Bay uses the northeastern coast's model and residual.}
\label{tab:a20_setup}
\begin{tabular}{lccccccc}
\hline
 & SD & \metric{params.err_pct_hi}th pct. & Above \metric{rev3.setup.threshold_m}~m & \multicolumn{2}{c}{Total SD} & Total SD & Largest bin \\
Sector & (m) & (m) & (\%) & (m) & (km) & over $Q$ & mean (m) \\
\hline
Kizlyar--Dagestan coast & \metric{rev3.setup.west.sd_m} & \metric{rev3.setup.west.abs_p95_m} & \metric{rev3.setup.west.share_above_0p20_pct} & \metric{rev3.setup.west.total_sd_m} & \metric{rev3.setup.west.total_sd_h_km} & \metric{rev3.setup.west.total_sd_over_floor} & \metric{rev3.setup.west.unit_mean_abs_max_m} \\
Volga--Ural interfluve coast & \metric{rev3.setup.interfluve.sd_m} & \metric{rev3.setup.interfluve.abs_p95_m} & \metric{rev3.setup.interfluve.share_above_0p20_pct} & \metric{rev3.setup.interfluve.total_sd_m} & \metric{rev3.setup.interfluve.total_sd_h_km} & \metric{rev3.setup.interfluve.total_sd_over_floor} & \metric{rev3.setup.interfluve.unit_mean_abs_max_m} \\
Northeastern coast & \metric{rev3.setup.northeast.sd_m} & \metric{rev3.setup.northeast.abs_p95_m} & \metric{rev3.setup.northeast.share_above_0p20_pct} & \metric{rev3.setup.northeast.total_sd_m} & \metric{rev3.setup.northeast.total_sd_h_km} & \metric{rev3.setup.northeast.total_sd_over_floor} & \metric{rev3.setup.northeast.unit_mean_abs_max_m} \\
Open shelf & \metric{rev3.setup.openshelf.sd_m} & \metric{rev3.setup.openshelf.abs_p95_m} & \metric{rev3.setup.openshelf.share_above_0p20_pct} & \metric{rev3.setup.openshelf.total_sd_m} & \metric{rev3.setup.openshelf.total_sd_h_km} & \metric{rev3.setup.openshelf.total_sd_over_floor} & \metric{rev3.setup.openshelf.unit_mean_abs_max_m} \\
Volga delta & \metric{rev3.setup.volga.sd_m} & \metric{rev3.setup.volga.abs_p95_m} & \metric{rev3.setup.volga.share_above_0p20_pct} & \metric{rev3.setup.volga.total_sd_m} & \metric{rev3.setup.volga.total_sd_h_km} & \metric{rev3.setup.volga.total_sd_over_floor} & \metric{rev3.setup.volga.unit_mean_abs_max_m} \\
Ural delta & \metric{rev3.setup.ural.sd_m} & \metric{rev3.setup.ural.abs_p95_m} & \metric{rev3.setup.ural.share_above_0p20_pct} & \metric{rev3.setup.ural.total_sd_m} & \metric{rev3.setup.ural.total_sd_h_km} & \metric{rev3.setup.ural.total_sd_over_floor} & \metric{rev3.setup.ural.unit_mean_abs_max_m} \\
Komsomolets Bay & \metric{rev3.setup.komsomolets.sd_m} & \metric{rev3.setup.komsomolets.abs_p95_m} & \metric{rev3.setup.komsomolets.share_above_0p20_pct} & \metric{rev3.setup.komsomolets.total_sd_m} & \metric{rev3.setup.komsomolets.total_sd_h_km} & \metric{rev3.setup.komsomolets.total_sd_over_floor} & \metric{rev3.setup.komsomolets.unit_mean_abs_max_m} \\
\hline
\end{tabular}
\end{table}
